% SPDX-FileCopyrightText: 2026 Will Cook
% SPDX-License-Identifier: CC-BY-4.0
% System-architecture paper; build with tectonic.
\documentclass[10pt]{article}
\usepackage{paper-house-style}
\usepackage{booktabs,array,float,placeins,multicol}
\usepackage{tikz}
\usetikzlibrary{arrows.meta,positioning,calc,fit,shapes.geometric}
\usepackage[colorlinks=true,linkcolor=linkinternal,urlcolor=linkexternal,citecolor=linkinternal]{hyperref}
\hypersetup{
  bookmarksnumbered=true,bookmarksopen=true,bookmarksopenlevel=1,
  pdftitle={Problem-Sized Lean Worlds},
  pdfauthor={Will Cook},
  pdfsubject={Architecture and trust boundaries of an AI-assisted Lean research system},
  pdflang={en-GB},
  pdfkeywords={formal mathematics, Lean 4, AI agents, research architecture, claim verification, reproducibility},
}
\newcommand{\repobase}{https://github.com/wcook04/plectis-erdos}
% BEGIN generated_semantic_coverage_macros
\newcommand{\SemanticAuthoredTheoremLike}{145,515}
\newcommand{\SemanticAuthoredInterpreted}{139,817}
\newcommand{\SemanticDirectEvidence}{3,329}
\newcommand{\SemanticContextual}{136,488}
\newcommand{\SemanticStructuralOnly}{5,698}
\newcommand{\SemanticAuthoredPercent}{96.1}
% END generated_semantic_coverage_macros
\newcommand{\repolink}[2]{\href{\repobase/blob/9b654f4cce44384f69d21fd2d0c51e62c4812b76/#1}{#2}}
\emergencystretch=1.7em

\runninghead{Problem-sized Lean worlds}
\title{Problem-Sized Lean Worlds}
\subtitle{An authority-separated architecture from AI search to public mathematical claims}
\author{Will Cook}
\date{September 2026}

\begin{document}
\maketitle
\paperprovenance
\fontsize{9.2pt}{11pt}\selectfont

\begin{abstract}
This paper describes an implemented architecture for AI-assisted mathematics
under proof abundance. It separates reasoning scope, permission to change
shared state, validation capacity, evidence class, public-claim authority,
and reviewer attention. Each transition requires evidence appropriate to the
next stage. Recorded failures can inform later search and validation; their evidential
value depends on what the attempt established.

Work is organised around \emph{problem-sized mathematical worlds}: formal
declarations, experiments, source literature, failed mechanisms, open
obligations, and a reviewed public boundary. Formal dependencies, authored
mathematical meaning, and public claims have separate graphs. In the private
workbench, a compiled comprehension packet gives an agent the endpoint and
claim ceiling before proposing a route. It federates corpora by content
digest, distinguishes proved results from open producers and closed routes,
and reports omissions and exceeded context budgets. The contribution is the
connection between research state, scoped changes, formal checks, and public
explanation, with their evidential roles kept distinct. The related-work
comparison claims no priority.

Computation can reject a route or supply finite evidence. Lean checks a proof
of the formal statement in the source; mathematical review must still
establish that the statement captures the intended problem and that the paper
describes it accurately. Comparator checks selected propositions and their
axiom boundary. Local review selection prepares result families for external
registration and review. These checks leave novelty, importance, acceptance,
and open-problem status to separate assessments.

The public repository supplies Lean source, a claim record, papers,
navigation, and release checks that can be inspected from a fresh clone. A
worked example follows a finite certificate for Erd\H{o}s Problem 249 through
the checks that prevent an unbounded public claim. The system is a working
prototype; multi-user use has not been validated. As of 31 August 2026, the
author had recorded neither a completed external cold-clone use nor an
accepted external contribution. All eight problems remain open.
\end{abstract}

\begin{paperresult}{Main contribution and claim ceiling}
\textbf{Contribution.} The prototype keeps six authorities separate and
makes every transition from conjecture to public wording carry its evidence
and its limit.  Problem-sized worlds join formal dependencies, mathematical
interpretation, failed routes, and reviewed claims without treating any one
graph as authority for the others.  \textbf{Demonstration.} The mechanism is
exercised on a self-contained public Lean corpus and an audited finite
certificate for Erd\H{o}s Problem 249.  \textbf{Ceiling.} No open problem is
claimed solved, and external use remains untested.
\end{paperresult}

\section{The problem: many kinds of evidence}

Suppose an AI system proposes a proof of a mathematical statement. Several
different questions immediately arise.

\begin{enumerate}
\item Did the system understand the intended problem?
\item Did it explore the important alternatives and notice counterexamples?
\item Does the proposed formal statement say what the informal statement says?
\item Does Lean accept a proof of that exact formal statement?
\item Does the public explanation stay within the proved scope?
\item Has anyone outside the producing system reviewed, accepted, or absorbed
      the result?
\end{enumerate}

These questions require different tests. A numerical experiment can expose a
false conjecture; Lean checks a formal theorem; a release checker preserves a
recorded decision about wording. None settles the other questions, and
external attention alone does not validate a proof. The architecture records
how work passes between these stages.

The \href{\repobase}{public repository studied here} contains Lean source,
papers, and release machinery around eight open Erd\H{o}s problems. At this
revision all eight problems remain open. The repository contains substantial
intermediate theorems, exact reformulations, conditional reductions, finite
certificates, and no-go results. The surrounding private workbench is broader:
it supports research, agent coordination, computational experiments,
formalisation, exposition, and controlled public projection. Private machinery
explains how the work was produced; it supplies no hidden proof authority to
the public clone.

The central engineering problem is preserving the meaning of a result when
it crosses these representations. A theorem name locates a declaration; it
does not specify which informal sentence that declaration supports. Conversely,
a paper can explain a sound argument whose steps have not all been formalised.
The implementation therefore records the correspondence explicitly and checks
that subsequent edits preserve it. Section~\ref{sec:example} follows one such
correspondence from a finite certificate to its public limitation. The
\repolink{docs/ARCHITECTURE.md}{public architecture guide} supplies the
clone-local map; the private mechanisms below describe the production design
and are not additional services that a reader must install.

\paragraph{Contribution and scope.}
The proposed contribution is an executable \emph{claim-transition
architecture}, assembled from established components including agents,
queues, file locks, theorem graphs, Lean checking, pull requests, and credit
records. Reasoning, mutation, validation, interpretation, publication, and
review have separate resource limits and permissions. Explicit fan-in gates
bring their outputs together. Proofs, counterexamples, no-gos, and unresolved
obligations pass through the same machinery with their evidence classes
preserved. A local failure can change a later agent's route, context, lease,
experiment, or validator without changing a mathematical statement's truth
status.

Four terms distinguish implementation from proposed use. \emph{Implemented}
means that source and an executable check or receipt exist. \emph{Projection}
is a generated view over more authoritative records. \emph{Inactive} means
that implemented machinery was not running at the reported snapshot.
\emph{Hypothesis} denotes a proposed experiment. The router, work leases,
corpus maps, Lean gates, Comparator, review-selection records, and release
checks are implemented; dashboards and graph views are projections. The
resident maintenance daemon was inactive at one recorded snapshot. Mass
frontier-model mining and training on the no-go graph remain hypotheses.

\subsection{Two architectural claims}

\paragraph{Problem neighbourhoods.}
An agent works in a bounded neighbourhood of a problem-sized world. The
neighbourhood starts with the exact endpoint and current claim ceiling, then
selects established premises, nearby declarations, open producers, consumers,
alternative coordinates, executable experiments, scoped falsifiers and
no-gos, source literature, and public boundaries. Each edge records its basis
and evidential limits; the packet lists omissions and expansion routes. File
proximity, lexical similarity, and model interpretation help navigation but
supply no proof.

\paragraph{Separate permissions and evidence.}
The architecture treats six resources as non-fungible. Reasoning scope does
not confer a write lease, and a write still requires validation. A Lean
receipt establishes formal acceptance, leaving semantic review, novelty, and
community acceptance to their respective gates. Typed artefacts record these
transitions.

\begin{table}[!ht]
\centering\scriptsize
\begin{tabular}{@{}>{\raggedright\arraybackslash}p{0.19\linewidth}
                    >{\raggedright\arraybackslash}p{0.30\linewidth}
                    >{\raggedright\arraybackslash}p{0.39\linewidth}@{}}
\toprule
\papertablehead
Resource & Governing object & What possession cannot imply \\
\midrule
Reasoning scope & Task-conditioned context and trace & Mutation permission or mathematical truth \\
Mutation permission & Exact path/work lease and change & Formal acceptance or public promotion \\
Validator capacity & Single-flight slot and build receipt & Theorem failure when a request is deferred \\
Evidence class & Experiment, counterexample, theorem, or authored relation & Automatic promotion into another class \\
Public-claim permission & Reviewed claim record and release relationship & Novelty, independent review, or acceptance \\
Reviewer attention & Ranked packet and recorded human outcome & Proof authority or canonical status \\
\bottomrule
\end{tabular}
\caption{The six non-fungible resources of the claim-transition architecture.}
\label{tab:nonfungible}
\end{table}

Mathematical and operational evidence update different graphs. The
mathematical graph records theorems, counterexamples, no-gos, and bounded
experiments with their evidence classes. The control graph records route
misses, stale views, repeated workarounds, validation failures, and resource
conflicts. After a generalisation guard, these may justify changes to a
router, skill, check, or standard. Reviewed claims and publication artefacts
connect the graphs without allowing either to determine the other's
conclusions.

\section{The whole lifecycle in one picture}\label{sec:lifecycle}
\papersectiontarget{systems-lifecycle}

Figure~\ref{fig:lifecycle} shows the main path and its feedback loop. Failed
proofs, reviewer objections, and public drift can inform earlier stages by
identifying errors or better methods. This feedback does not alter what Lean
checked or make an unproved statement true.

\begin{figure}[!t]
\centering
\begin{tikzpicture}[
  box/.style={paper stage, text width=2.05cm},
  private/.style={box, fill=MidnightBlue!5, draw=MidnightBlue!55},
  formal/.style={box, fill=TealBlue!6, draw=TealBlue!60!black},
  public/.style={box, fill=ForestGreen!6, draw=ForestGreen!55!black},
  assure/.style={box, fill=BurntOrange!7, draw=BurntOrange!65!black},
  arr/.style={paper arrow},
  back/.style={paper feedback},
  lab/.style={font=\scriptsize\sffamily\bfseries, text=gray!70!black}
]
\node[private] (intent) at (0,0) {{\scriptsize\sffamily\bfseries 1 \; SELECT}\\[3pt]
  \textbf{Intent and route}\\[2pt]\scriptsize question, scope, claim};
\node[private, right=3.5mm of intent] (reason) {{\scriptsize\sffamily\bfseries 2 \; TEST}\\[3pt]
  \textbf{Reason and test}\\[2pt]\scriptsize candidates and failures};
\node[formal, right=3.5mm of reason] (lean) {{\scriptsize\sffamily\bfseries 3 \; CHECK}\\[3pt]
  \textbf{Formalise}\\[2pt]\scriptsize proof and kernel verdict};
\node[public, right=3.5mm of lean] (claim) {{\scriptsize\sffamily\bfseries 4 \; EXPLAIN}\\[3pt]
  \textbf{Public claim}\\[2pt]\scriptsize result, evidence, limit};
\node[assure, right=3.5mm of claim] (review) {{\scriptsize\sffamily\bfseries 5 \; RELEASE}\\[3pt]
  \textbf{Assure}\\[2pt]\scriptsize replay and boundary checks};
\draw[arr] (intent) -- (reason);
\draw[arr] (reason) -- (lean);
\draw[arr] (lean) -- (claim);
\draw[arr] (claim) -- (review);
\draw[back] (review.south) to[out=-115,in=-65,looseness=1.15]
  node[lab, below=4pt] {failures and review return as evidence; authority stays put} (reason.south);
\node[font=\scriptsize\sffamily\bfseries, text=MidnightBlue!75!black,
  above=5mm of $(intent.north)!0.5!(reason.north)$] {PRIVATE WORKBENCH};
\node[font=\scriptsize\sffamily\bfseries, text=ForestGreen!62!black,
  above=5mm of $(lean.north)!0.5!(review.north)$] {SELF-CONTAINED PUBLIC ARTEFACT};
\draw[paper boundary] ($(reason.east)!0.5!(lean.west)+(0,-9mm)$) --
  ($(reason.east)!0.5!(lean.west)+(0,12mm)$);
\end{tikzpicture}
\caption{The end-to-end architecture. The dashed vertical line marks the
release boundary. Everything to its right remains usable without the private
workbench. The labels describe the production deployment; a contributor can
also select and test new mathematics using the public clone.}
\label{fig:lifecycle}
\end{figure}

The lifecycle has four recurring operations. First, \emph{selection}: choose a
bounded object rather than treating the whole repository as one prompt.
Second, \emph{execution}: run an experiment, edit a proof, or write an
explanation. Third, \emph{validation}: ask the authority appropriate to that
object. Fourth, \emph{binding}: record the result, limitation, and source so a
later agent does not have to rediscover them. The final feedback step turns a
local success or failure into a reusable route, check, or warning when it
genuinely generalises.

The public \repolink{scripts/proof_cockpit.py}{proof-cockpit card}
gives an agent this separation.
It reads checkout and toolchain
identity, corpus scale, registered open propositions, problem obligations, and
workbench sessions from public files.  It imports no private state and adds no
evidence layer to Figure~\ref{fig:lifecycle}.  Its \texttt{--check} mode
runs the claim-registry, cold-clone, and projection-freshness checks; neither
mode invokes Lean or decides whether prose follows from a theorem.

Once a person has compared a formal theorem with its public wording, a program
can preserve the resulting decision about names, files, wording, and limits.
It cannot decide whether the decision was correct.  The workflow does not
technically force a second independent mathematician: one maintainer can edit
the Lean statement, record, and prose together so that every comparison agrees
with the same mistake.
\subsection{Skill-addressable job lifecycles}
\papersectiontarget{systems-job-lifecycle}

The public clone gives each recurring job one entry skill. A skill owns the
job's starting state, permitted mutations, required evidence, stopping or
re-entry condition, and next owner; it does not acquire the authority of the
objects it coordinates. Figure~\ref{fig:job-lifecycle} records the ordinary
research path.

\begin{figure}[!t]
\centering
\begin{tikzpicture}[
  job/.style={paper stage, text width=2.35cm, minimum height=1.25cm},
  arr/.style={paper arrow}
]
\node[job, fill=MidnightBlue!5] (invoke) at (0,0)
  {{\scriptsize\sffamily\bfseries 1}\\[2pt]\textbf{Invoke}\\[2pt]\scriptsize bounded route};
\node[job, fill=MidnightBlue!5, right=5mm of invoke] (work)
  {{\scriptsize\sffamily\bfseries 2}\\[2pt]\textbf{Work}\\[2pt]\scriptsize continue or suspend};
\node[job, fill=TealBlue!5, right=5mm of work] (validate)
  {{\scriptsize\sffamily\bfseries 3}\\[2pt]\textbf{Validate}\\[2pt]\scriptsize computation or Lean};
\node[job, fill=TealBlue!5, right=5mm of validate] (propagate)
  {{\scriptsize\sffamily\bfseries 4}\\[2pt]\textbf{Propagate}\\[2pt]\scriptsize classify consequences};
\node[job, fill=ForestGreen!6, below=8mm of propagate] (package)
  {{\scriptsize\sffamily\bfseries 5}\\[2pt]\textbf{Package}\\[2pt]\scriptsize evidence and provenance};
\node[job, fill=ForestGreen!6, left=5mm of package] (propose)
  {{\scriptsize\sffamily\bfseries 6}\\[2pt]\textbf{Propose}\\[2pt]\scriptsize pull request or issue};
\node[job, fill=BurntOrange!7, left=5mm of propose] (adopt)
  {{\scriptsize\sffamily\bfseries 7}\\[2pt]\textbf{Review and adopt}\\[2pt]\scriptsize reconcile, accept, credit};
\draw[arr] (invoke) -- (work);
\draw[arr] (work) -- (validate);
\draw[arr] (validate) -- (propagate);
\draw[arr] (propagate) -- (package);
\draw[arr] (package) -- (propose);
\draw[arr] (propose) -- (adopt);
\end{tikzpicture}
\caption{The clone-local job lifecycle. Closure means that the present delta
has evidence and dispositions; it does not mean that the open problem is
solved.}
\label{fig:job-lifecycle}
\end{figure}

The stable roles are orientation, bounded discovery, Lean validation,
consequence propagation, return packaging, proposal, and maintainer adoption.
Problem admission and manuscript revision have separate lifecycles. The
\repolink{skills/README.md}{skill catalogue} lists the current entry names and
links to definitions containing the executable details.

\section{The private workbench}\label{sec:private}

\subsection{Disk is shared memory}

Durable project state lives in source files, structured records, append-only
events, generated views, validation receipts, and authored explanations.
Conversations propose changes; files and their governing checkers determine
whether those changes exist. To resume work, the next agent reads a bounded
current-state packet and the relevant sources, so it need not rely on a
summary of an earlier conversation.

File classes have distinct roles. Raw operator language preserves the
request, and work records identify the claimed work and its owner. Source
files contain implementations and proofs; receipts record what ran. Authored
papers explain the work, while generated indexes provide navigation back to
the governing records.

At entry, a deterministic router selects a small context packet containing
candidate objects, authority boundaries, relevant standards, and permitted
next actions. The packet guides source inspection; its summary has no proof
authority.

\subsection{Type A and Type B}

\emph{Type A} means a coding agent working in the repository, such as Claude
Code, Codex, Cursor, Antigravity, or OpenCode. \emph{Type B} means a chat AI
used through a web interface: the operator gives it a selected research packet
and brings its response back to the repository. This is a distinction in
substrate access, not model quality: it does not rank either AI's intelligence.

\begin{table}[!ht]
\centering\small
\begin{tabular}{@{}>{\raggedright\arraybackslash}p{0.16\linewidth}
                    >{\raggedright\arraybackslash}p{0.34\linewidth}
                    >{\raggedright\arraybackslash}p{0.39\linewidth}@{}}
\toprule
\papertablehead
Actor & What it can do & Governing limit \\
\midrule
Type A & Inspect live project state; use repository tools; edit claimed files;
run tests; bind receipts and status. & It may change only the substrate and
paths its task authorises, and its claims remain limited by the relevant
validator or human review. \\
\addlinespace
Type B & Read a selected packet in web chat; return research, critique, or a
candidate for the operator to carry back. & It has
no direct access to these project files or tools. A Type A agent or the operator must
check and apply useful output. \\
\bottomrule
\end{tabular}
\caption{Type A/B is an authority distinction. Either type may be weak or
strong; a delegated tool-using agent is still Type A if it has live substrate
access.}
\label{tab:typeab}
\end{table}

Type B output can inform research without being treated as an observation of
private state. Type A changes remain attached to their paths, tests, and
receipts. Delegation alone does not make a worker Type B: its access to the
repository determines the role.

\subsection{Concurrency without shared-state confusion}

Several agents may work at once, but concurrency is admitted only where the
write scopes can be separated. A work item identifies the objective and
expected evidence. Before mutation, an agent claims exact paths for a bounded
lease. Independent claims may proceed concurrently; overlapping claims must be
coordinated or deferred. Fan-out is followed by a fan-in barrier where the
results are compared, validated, and integrated.

Append-only ledgers and immutable receipts record coordination state, and
generated status pages present views over them. They distinguish an active
agent from an existing change and an accepted change. Bounded job counts and
focused Lean builds limit resource use and prevent concurrent builds from
corrupting one another's evidence.

\subsection{The control plane is executable}

The router exposes typed entries for skills, standards, paper modules, live
work, and source artefacts. It selects a task-conditioned packet before the
agent opens large sources. A cold agent can proceed from a one-line flag to
a short card, a working context, and finally the governing file. Each layer
identifies its source, expansion route, and authority, making the selection
of context inspectable.

The work ledger records completed, reopened, and superseded work as
append-only lifecycle events. Present permission comes from short leases over
exact paths or work objects. Directory claims collide with their children;
expired leases lose authority unconditionally. A crashed holder leaves an
expiry and dirty-handoff record. Generated cohort views report activity,
collisions, and unknown scope, but cannot grant permission to write.

A resident-runtime layer called metabolism turns repository events into
bounded maintenance jobs. One daemon owns a SQLite store in write-ahead-log
mode containing events, jobs, runs, provider budgets, heartbeats, and temporal
blackboard claims. Hooks, filesystem scans, provider interruptions, and
explicit commands feed the store. Stable digests deduplicate events; an
allowlist limits execution; cooldowns regulate provider use. Expired owners
are recoverable, and a single-resident guard prevents competing schedulers.
Status pages report the store's state without authorising repairs. At one
snapshot used for this revision, the status surface reported the daemon as
not running while queued jobs remained visible. That snapshot demonstrates
persistent state in the implemented machinery, not an active maintenance
service.

The operating loop observes the current state, classifies the task, routes it
to an owner, claims the work, acts, validates, records the result, propagates
reusable lessons, and observes again. Continuous work repeats this loop with
explicit checkpoints and a new receipt for each advance.

\subsection{What continuous mathematical work looks like}

A continuous goal records a fixed research question, a changing frontier, and
an explicit claim ceiling. Each run resumes from the latest accepted state:
sources read, routes ruled out, computations interpreted, formal obligations
still open, and the precise statement that remains unproved. The agent
selects the highest-value available transition, such as closing a lemma,
producing a counterexample, sharpening a reduction, repairing a source claim,
or isolating a smaller obstruction.

The public continuation interface makes part of this design inspectable.
\repolink{scripts/continue_research.py}{The continuation program}
records a bounded session at a named starting commit, checks a returned result
against that session, and packages the return for inspection. It composes
existing query, workbench, and return-validation tools. It does not execute
commands merely because they appear in a contributor's return. Its optional
replay checks stored session probes; it is not a repository build or release
check. Persistence here means that the next contributor can recover the
question, starting source, attempted work, and evidence without recovering the
original conversation.

A compressed trace can omit paths or worker outcomes. Repository state,
command results, diffs, formal builds, and commit objects must establish
whether a change occurred. The workbench therefore records completeness
explicitly. Private activity traces document production; their event counts
are not used here as public evidence of throughput or mathematical progress.

Search and validation are interleaved, with scoped failed routes retained and
concurrent work permitted where tasks can be separated. A claimed advance
requires the relevant artefact and receipt. If an unavailable validator, a
human decision, or a new idea is needed, the run preserves its endpoint and a
re-entry point. During long computations, an agent can advance independent
proof, source, exposition, or adversarial work where it is safe to do so.

\subsection{Coupled continuous goals: discovery and stewardship}
\papersectiontarget{systems-coupled-goals}

Discovery and stewardship address different scales of the corpus. A
problem-directed miner works near the proof frontier; assessing whether its
new lemma is routine, weaker than an older theorem, or prominent enough to
lead a paper requires a wider view. The \emph{discovery goal} reads the current
problem world, runs discriminating computations, attempts proofs, records
falsifiers, and returns the smallest stable mathematical delta. The
\emph{stewardship goal} compares that delta with the current corpus, groups
related declarations into result families, and reconciles downstream uses.
The two roles retain separate responsibilities and authority.

\begin{figure}[!t]
\centering
\begin{tikzpicture}[
  role/.style={paper stage, text width=3.05cm, minimum height=1.5cm},
  object/.style={role, fill=TealBlue!6, draw=TealBlue!60!black,
    text width=2.72cm},
  output/.style={role, fill=ForestGreen!6, draw=ForestGreen!55!black,
    text width=3.15cm},
  arr/.style={paper arrow},
  feedback/.style={paper feedback}
]
\node[role, fill=MidnightBlue!5, draw=MidnightBlue!55] (mine) at (0,0)
  {{\scriptsize\sffamily\bfseries DISCOVERY}\\[3pt]\textbf{Test the frontier}\\[2pt]
   \scriptsize compute, prove, or falsify};
\node[object, right=7mm of mine] (delta)
  {{\scriptsize\sffamily\bfseries HANDOFF}\\[3pt]\textbf{Landed delta}\\[2pt]
   \scriptsize theorem, no-go, computation};
\node[role, fill=BurntOrange!7, draw=BurntOrange!65!black,
  right=7mm of delta] (steward)
  {{\scriptsize\sffamily\bfseries STEWARDSHIP}\\[3pt]\textbf{Place the result}\\[2pt]
   \scriptsize appraise, reconcile, explain};
\node[output, below=9mm of delta] (surfaces)
  {{\scriptsize\sffamily\bfseries PUBLIC SURFACE}\\[3pt]
   \textbf{Current mathematical picture}\\[2pt]
   \scriptsize paper order, evidence, frontier};
\draw[arr] (mine) -- (delta);
\draw[arr] (delta) -- (steward);
\draw[arr] (steward) -- (surfaces);
\draw[feedback] (surfaces.west) to[out=180,in=-90]
  node[below left, font=\tiny\sffamily, align=center]
  {ranked next question, missing evidence,\\discriminating test} (mine.south);
\end{tikzpicture}
\caption{The coupled continuous-goal protocol. A landed mathematical delta
is a reason to resume stewardship; a changed appraisal or consumer gap can
change the next mining target. The protocol calls for a new pass only after
a relevant change.}
\label{fig:coupled-goals}
\end{figure}

The stewardship pass makes four decisions separately. \emph{Authority}
records what Lean, a computation, a paper argument, or an external source
actually establishes. \emph{Mathematical appraisal} asks about logical reach,
mechanism depth, independence, sharpness, reuse, and the surviving open
boundary. \emph{Exposition placement} determines what leads the abstract,
which theorem receives the longest explanation, and what remains subordinate.
\emph{Work allocation} determines which missing implication or consumer
deserves the next unit of compute or expert attention. A result can be highly
significant but mechanically unready, fully checked but mathematically
routine, or ideal as a worked example without being the corpus's strongest
theorem. No single scalar or status is allowed to decide all four questions.

The steward can detect that a purported advance restates the open problem,
that several theorem names express one result, or that a paper buries a
stronger theorem. It can also find weaker Comparator interfaces and Palomar
packets that omit the hard mechanism. It may narrow, regroup, reorder, or
defer these presentations. Such judgements do not supply a second proof
authority: Lean validation, intended-meaning review, prior-art assessment, and
community acceptance remain separate.

The clone-local
\repolink{skills/run-coupled-research-goals/SKILL.md}{coupled-goals skill}
specifies the division through
\repolink{skills/mine-open-problem/SKILL.md}{bounded discovery} and
\repolink{skills/propagate-research-consequences/SKILL.md}{consequence propagation},
with source-pinned hand-offs between the jobs. The skill instructs an agent or
controller; it supplies no resident scheduler. One agent can perform the jobs
sequentially, or separate tasks, machines, models, or people can hold the roles.

A landed theorem, counterexample, authority change, paper correction, or
external review outcome triggers a stewardship pass. The pass produces a
source-pinned appraisal and a changed consumer or frontier disposition. This
can direct discovery towards a stronger target, a missing hypothesis, a
counterexample request, or a cheap discriminating computation. When neither
the frontier nor a downstream consumer has changed, both goals yield.

Before either role begins, \texttt{explain-public-system} provides a separate
cold-clone orientation job. An agent can read the public corpus and companion
papers, adapt the explanation to a lay reader, mathematician, formaliser,
compute contributor, reviewer, or infrastructure contributor, and point to the
exact evidence behind its account. The human therefore need not learn the file
layout before asking a useful question. Explanation remains a projection: it
cannot promote its own summary into proof or acceptance.

\section{The mathematical reasoning loop}\label{sec:mathloop}
\papersectiontarget{systems-mathloop}

Mathematical work starts by stating the problem in ordinary language with its
known status and source. Candidate mechanisms are ranked by endpoint
proximity, depth, independence from other routes, usefulness for future work,
and overclaim risk. Theorem count and ease of formalisation do not determine
that ranking. A decisive obstruction or close conditional reduction may
deserve more attention than many routine lemmas.

\subsection{Experiments are route selectors}

Computation serves three useful roles. It can find a counterexample, measure a
finite pattern, or test whether a proposed mechanism is plausible enough to
formalise. Each experiment records its inputs, code, finite domain, output, and
interpretation. The interpretation must state what the experiment cannot show.

A finite computation becomes formal evidence only through an explicit bridge.
For example, Lean may evaluate a finite proposition with \texttt{decide} and
check the resulting proof term. Even then, the conclusion remains finite. A
pattern observed for many inputs does not become an ``all inputs'' theorem, and
a successful numerical approximation does not become an exact equality. Failed
experiments are kept when they prune a natural route; otherwise later agents
pay to repeat the same mistake.

\paragraph{Negative evidence.}
A failed agent attempt records an uncompleted search path. A counterexample
refutes a conjecture on its stated domain, while a Lean no-go theorem rules
out a strategy class under explicit hypotheses. Counterexamples and no-gos
can prevent repeated work and identify a missing ingredient for the endpoint.
The corpus keeps problem-specific reasoning, proof gaps, coefficient-only
and fixed-precision obstructions, and corrected statements alongside proved
theorems. Each no-go retains its scope, so the failure of one mechanism does
not rule out other proofs.

\subsection{Status, formal acceptance, and mathematical judgement}

The workflow distinguishes three senses of ``verified''. Experiments guide
route selection without supplying any of these judgements by themselves.

\begin{enumerate}
\item A \emph{status oracle} says what a public registry or the literature
      currently reports about the problem.
\item A \emph{formal oracle} says whether the pinned Lean environment accepts
      this exact statement and proof under the permitted axioms.
\item A \emph{research-validity judgement} asks whether the formal statement
      matches the intended problem and whether the result matters in context.
\end{enumerate}

The second can be mechanical. The first can become stale. The third remains a
mathematical and scholarly judgement. A successful Lean build without statement
reconciliation is therefore not the end of the reasoning process.

Each iteration begins with the existing Lean and semantic graph. Proof search
and bounded computation probe gaps and patterns that may be too distributed
for an unaided manual scan. A checking pass records the kernel result,
reconciles formal and informal statements, and packages a theorem,
counterexample, or diagnosed failure. Accepted nodes and edges inform the
next search; whether they improve it remains an evaluation question. Reusable
process lessons can change heuristics and context, without changing
mathematical truth. For mechanisms drawn from the literature, annex records,
locators, and public citations retain the original authorship and claim scope.
Lean checks the new formal statement.

\subsection{From local progress to reusable mathematics}

A lemma found in one Erd\H{o}s problem may yield a comparison principle,
finite combinatorial device, analytic estimate, or Lean interface useful in a
wider library. The system distinguishes \emph{process up-propagation}, which
improves how later agents work, from \emph{mathematical canonicalisation},
which develops a local result into a reusable theorem.

Canonicalisation requires examining which hypotheses the proof uses, removing
problem-specific coordinates, checking prior art and existing library
interfaces, and stating a natural general theorem. The local result should
be an explicit specialisation. Dropped hypotheses require counterexample
tests, and the general proof needs its own exact check. At least one additional
consumer or explanatory use should establish the value of the generalisation;
broader notation alone is insufficient.

This yields three statuses that must not be collapsed: checked inside one
problem world; proposed as reusable mathematics; and suitable for an external
shared library. The present system can preserve the first, help construct and
test the second, and prepare evidence for the third. It cannot grant the third
status to itself. Mathlib is open to new contributors, but its current guide
sets high standards for generality, integration, maintainability, style,
documentation, and responsible disclosed use of AI \cite{mathlib}. A subject
and Lean expert must decide whether a candidate belongs there, reshape and
review it as needed, and take responsibility for any upstream discussion or
pull request. The originating route should remain visible in provenance; the
expert should receive explicit credit for the generalisation, library design,
formalisation, review, and stewardship actually supplied. The
\papersectionlink{open-source-mathematics-strategy.pdf}%
{strategy-local-to-general}{open-source strategy} gives the corresponding
participation and credit protocol.

\subsection{Problem-sized Lean worlds and bounded theorem neighbourhoods}

The public corpus provides several maps. The
\repolink{docs/declaration_atlas.json}{declaration atlas} locates statements.
\repolink{scripts/query_corpus.py}{Dependency queries} distinguish module
imports from constant references extracted from elaborated Lean declarations;
unavailable or stale extraction cannot establish that a dependency is absent. The
\repolink{docs/semantic/README.md}{authored semantic layer} explains
mathematical roles such as specialisation, transport, and scoped obstruction.
The \repolink{docs/claims.json}{claim record} connects selected results to
reviewed wording and exact open propositions. A dependency can guide premise
selection without establishing an authored analogy, and an analogy can
suggest a proof without supplying one.

The public Problem 1041 materials illustrate why these distinctions matter.
The \repolink{docs/problems.json}{problem index} identifies the integrated
library and its separately governed research export. The export's
\repolink{research_corpus/Erdos1041/CORPUS_MANIFEST.json}{manifest} inventories
its files; its
\repolink{research_corpus/Erdos1041/FRONTIER.md}{frontier account} records
surviving mechanisms and refutations; and its
\repolink{research_corpus/Erdos1041/STRONGEST_RESULTS.json}{selected-results map}
binds statements to source with their premises and falsifiers. An exported
file, an integrated module, a semantic interpretation, and a reviewed public
claim are different objects. The declaration and semantic inventories expose
their respective coverage; generated certificate families make raw theorem
counts a poor measure of mathematical value.

The problem cockpit first fixes the endpoint, status, claim ceiling, and
canonical frontier. It distinguishes open producers from the target, so a
promising route does not replace the problem. A bounded neighbourhood then
selects the strongest relevant results, exact source coordinates, imports, consumers,
alternative formulations, experiments, counterexamples, no-gos, and
literature. Each item has an evidence class. An omission receipt records what
was excluded and how to retrieve it. For an exact claimed path, a local
connection card can present prerequisites, sibling mechanisms, consumers,
falsifiers, and the validation target before broader retrieval.

\subsection{Comprehension before the mathematics}
\papersectiontarget{systems-comprehension}

The private workbench compiles this neighbourhood into a packet conditioned on
the agent's query. The following federation and inference-workspace behaviour
belongs to that workbench. In a public clone, the independently usable
\repolink{scripts/query_corpus.py}{corpus query tool} supplies local claim,
declaration, dependency, and connection-card routes; it does not require the
private compiler or its attached corpora.

Infrastructure questions have a separate entry path. The private workbench
offers a small choice of architecture, source layout, local connections,
validation, experiment evidence, and paper ownership; the agent selects an
existing owner without scanning the proof corpus to discover its organisation.
The public clone's \repolink{scripts/agent_entry.py}{task-entry tool} routes
architecture and navigation work to its
\repolink{skills/maintain-public-infrastructure/SKILL.md}{infrastructure
maintenance workflow}. In both settings, source topology locates work, module
imports guide inspection, and a build plan selects validation. These routes
help an agent choose what to inspect or run; passing their structural checks
does not establish a mathematical claim.

Each corpus retains its own authority. The compiler stores compact
descriptors and content fingerprints for the private projection and attached
public checkouts, without copying them into a central index. It opens an
attached exhaustive atlas only after the identities reconcile and bounded
retrieval fails to answer the query. It does not rebuild the expensive Lean
microcosm. Every packet names the corpora consulted and their digests.

The packet leads with the endpoint. A problem cockpit fixes the target
statement, its status, the claim ceiling, and the current claim frontier,
labelling each frontier claim with an evidence class and a logical altitude
that separates an equivalence from a reduction, a necessary condition, a finite
exclusion, and a no-go. Beneath it a mechanism landscape keeps proved or
kernel-checked results apart from open producers and from routes a
counterexample has closed. Retrieval moves between global frontier structure,
concept neighbourhoods, and premise or obligation ancestry, and exact imports,
authored arguments, generated navigation, lexical references, and semantic
inference remain visibly different edge classes.

An inference workspace organises this material in three planes. The exact
commitment plane holds source coordinates and statements. The strategy plane
branches over coherent premise groups, and the adversary plane records the
falsifiers each branch must survive. A changed corpus fingerprint invalidates
the commitments. Every packet states the limit of this planning workspace:
only Lean, an exact counterexample, or an owner verifier receipt may change
a claim's status.

The compiler reports a query as unanchored if it names no theorem,
declaration, or claim. Its consumer must name one before treating a route as
target-specific. A packet that exceeds its requested context budget exits
non-zero, reports requested and estimated token counts, lists dropped planes,
and gives commands to restore them. If the adversary plane was omitted, the
packet records that omission and an expansion route; it does not report an
absence of falsifiers.

A correctly identified and well-formed packet can still select a weak
mechanism, omit the governing obstruction, or carry a mathematically mistaken
interpretation. The supported claim concerns its starting state:
bounded context with source digests, evidence classes, and declared
omissions. The packet has no proof authority. This paper does not measure
whether it helps agents produce better mathematics.

\subsection{Consequence propagation}

When a declaration lands, the direction reverses. A consequence mapper begins
at the exact changed object and enumerates reverse imports, claim records,
validators, experiments, papers, and open obligations that may now be stale.
A semantic second pass must choose: update now, verify unchanged, defer with a
reason, or mark outside scope. Empty lexical search is not evidence of no
consequence, and no projection may bulk-strengthen a family of claims. A formal
evidence cell carries the result, its explicit non-claim, evidence class,
source and receipt links, and next unresolved obligation through this fan-out.
The public
\repolink{skills/propagate-research-consequences/SKILL.md}{consequence-propagation skill}
specifies the same
discipline without depending on the private workbench. For work returned from
an older clone, it runs once against the contributor's recorded starting
commit and again after the reviewed change has been reconciled with current
main. The original delta and any conflict-resolution delta remain separate
evidence and receive separate attribution.

This same architecture distinguishes mathematical propagation from process
propagation. A theorem, counterexample, or no-go changes the mathematical graph
only through its own verifier. A reusable failure in navigation, scheduling,
experimentation, or validation enters a guarded packet containing the local
case, failure class, evidence, sibling scan, proposed owner, overgeneralisation
guard, validation, and stop condition. If accepted, it may change a route,
skill, check, or standard for later agents. The mechanism is implemented and
has worked examples, but there is not yet a complete measure of what fraction
of useful failures propagate. The supported claim is that failures \emph{can}
alter the durable control plane without altering mathematical truth.

\section{The public Lean repository}\label{sec:public}
\papersectiontarget{systems-public}

The public checkout contains the source needed to inspect its claims and
replay its checks, with no calls to the private workbench or unpublished
private lemmas. Builds also require the public toolchain and dependency
versions recorded by \repolink{lean-toolchain}{lean-toolchain} and
\repolink{lake-manifest.json}{lake-manifest.json}, which a fresh checkout may
need to download. Self-contained therefore describes the source and its
authority, rather than an offline installation. The private workflow remains
production provenance.

The repository has two Lean roots. The reviewed root contains the established
249/257 publication lane. The problem-owned expansion root contains work on six
additional Erd\H{o}s problems and unpromoted lanes for 249 and 257. Both roots
are checked by the same Lean kernel, but kernel acceptance does not automatically
promote a declaration into a reviewed public claim. Promotion is a separate
change to the claim record and exposition.

The public sources divide responsibility as follows:

\begin{table}[H]
\centering\small
\begin{tabular}{@{}>{\raggedright\arraybackslash}p{0.23\linewidth}
                    >{\raggedright\arraybackslash}p{0.31\linewidth}
                    >{\raggedright\arraybackslash}p{0.34\linewidth}@{}}
\toprule
\papertablehead
Surface & Authority & It does not establish \\
\midrule
Lean source and pinned toolchain & Exact formal statements and proofs accepted
by the kernel. & Intended meaning, novelty, significance, or faithful prose. \\
\addlinespace
Reviewed claim record & Approved wording, status, evidence, bounded domain,
and adjacent open statement for selected claims. & Correctness or completeness
of the human review. \\
\addlinespace
Papers and guides & Human explanation and reading order within the recorded
claim ceiling. & New formal authority. \\
\addlinespace
Generated maps and query tools & Bounded navigation across declarations,
problems, graphs, papers, and claims. & Proof or permission to strengthen a
claim. \\
\addlinespace
Release checks and continuous integration & That configured identities,
relationships, generated files, licences, and negative tests pass on a named
revision. & Understanding of unrestricted prose or independent mathematical
approval. \\
\bottomrule
\end{tabular}
\caption{The public authority split.}
\label{tab:authority}
\end{table}

The companion public Plectis repository has a different role. It publishes
runnable, bounded mechanism slices and receipts from the wider system. The Lean
repository publishes a mathematical corpus. Neither public repository inherits
authority from the other, and neither is a public mirror of the private root.

\section{Comparator, Palomar, and publication}\label{sec:assurance}

\subsection{Comparator: selected statement interfaces}

Comparator requires selected propositions to be stated again in a challenge
module without their proofs. A solution module must provide terms of those
exact types. The \repolink{verification/comparator.json}{configuration}
names the permitted axioms, modules, and runtime receipt; the
\repolink{docs/EXTERNAL_VERIFICATION.md}{verification guide} explains their
roles. A named altered statement must fail, testing that the harness is not
vacuous.

The check establishes that the proof-bearing corpus implements the separately
stated Lean interface under the specified axiom budget. It leaves the
translation from informal mathematics, novelty, importance, and independent
human review unresolved. The appropriate description is
``Comparator-checked''; the receipt alone does not establish independent
verification.

\subsection{Review selection and the Palomar registry}

Comparator inventories evidence. Local review selection groups declarations
into result families and ranks them by mathematical signal to choose a small
portfolio for review. Each packet keeps the hard mechanism beside the
surviving boundary. The
\repolink{docs/verification/PALOMAR_QUALIFICATION.md}{qualification record}
records whether the packet satisfies its structural requirements. The
\repolink{docs/PALOMAR_RESULT_SHOWCASE.json}{local selection record}
uses Palomar's name because selected families may be submitted to its external
registry of Lean-verified mathematics. Palomar documents mechanical proof
checks and automated editorial filtering, without human peer review or
significance ranking. Local qualification, registry inclusion, and human
mathematical review are separate events; none alone establishes broad
community acceptance or significance.

The local ranking asks which result changes the mathematical picture, which
conditional route approaches an endpoint, which obstruction saves repeated
work, and which explanation would help an expert assess the claim. These
criteria direct attention without altering proof status or rewarding theorem
count and generated volume by themselves.

\subsection{Publication and propagation}

After formal proof and statement reconciliation, a result may enter an authored
paper, a claim record, a Comparator packet, and a local review unit. Each
records a different decision with a different ceiling. The
\repolink{docs/publication_contract.json}{publication contract} inventories
the manuscripts and rendered files. The
\repolink{scripts/check_release.py}{release program} and
\repolink{.github/workflows/lean.yml}{continuous-integration workflow}
check the recorded relationships. Release policy requires both the Lean build
and the public-surface checks; the Python release program does not itself
elaborate Lean.

Failures and review feed back into the appropriate owner. A proof failure can
change a formalisation heuristic, a counterexample can close a route, a
reviewer objection can narrow a claim, and release drift can prompt a new
check. The resulting change belongs in a theorem, experiment record, skill,
standard, route, or claim boundary. It must not rewrite raw intent, generated
views, or mathematical status by implication.

\subsection{From formal checking to mathematical use}

Tao argues that AI mathematics should be judged along a chain: proof
generation, verification, exposition, publication and community digestion,
then eventual canonicalisation \cite{taoai}. This architecture is a concrete
attempt to build for that whole chain. The private reasoning and experiment
loops support generation. Lean checks formal proofs, while Comparator protects
selected exact interfaces. The problem papers, reasoning surfaces, graphs, and
claim records support exposition. Local review selection triages scarce expert attention, and
the contribution and release paths support attributable publication. Digestion,
acceptance, and canonicalisation remain achievements of the mathematical
community, never local status fields.

Tao warns that polished AI exposition can erase the friction that reveals
where the mathematical difficulty lies. The reasoning surfaces and no-go graph
retain false starts, corrected claims, missing producers, and formal
obstructions alongside the successful theorems. They help a reader understand
why the surviving boundary is difficult and where a new idea is needed.

Paper authoring itself participates in this loop. As the mathematical and
control graphs change, agents assemble and revise problem papers, architecture
papers, reviewer cards, and claim records from the source-current evidence.
We call this AI-assisted preparation \emph{pre-digestion}. Agents organise
arguments, expand intermediate steps, produce examples and diagrams, and link
statements to their sources. A missing explanation, unstable term, hidden
quantifier change, or unregistered limitation found during that work becomes
a defect to investigate in the theorem, claim record, or software.

A digestive-system analogy distinguishes breaking food down from absorbing
its nutrients. Similarly, preparing an argument in an accessible form does
not establish that a mathematician has absorbed its ideas. Mathematical
digestion requires people to reconstruct the reasoning, examine where the
assumptions matter, explain the difficult step, and connect the result to
mathematics they can use. AI can assist with this work; generated exposition
cannot certify that it has happened. This is our analogy for Tao's distinction
between verified proofs, readable exposition, and results understood and
accepted by mathematicians \cite[Sec.~6, pp.~6--8; Sec.~8, p.~11]{taoai}.
Formal checking, demonstrated human understanding, and community acceptance
remain separate forms of evidence.

This manuscript is a reflexive example: it was rewritten by the system while
agents inspected the mechanisms it describes, then compiled and checked as a
public artefact. Self-authorship records production and supplies no independent
validation; a clear account does not constitute an independent review.

Tool use is disclosed and models are not listed as authors. Human
contributors receive differentiated credit, and references and source paths
remain inspectable. Automatic checks filter work without replacing peer
review; theorem counts do not determine mathematical value. These choices
make the repository an implementation experiment in preparing a small
project for proof abundance.

\section{One complete boundary: finite is not unbounded}\label{sec:example}

Erd\H{o}s Problem 249 asks whether
\[
S=\sum_{n\ge1}\frac{\varphi(n)}{2^n}
\]
is irrational, where \(\varphi(n)\) is Euler's totient function
\cite{erdosgraham}. We use this example because its finite and unbounded
statements can be followed through every publication layer; it is not a
ranking of the corpus's strongest mathematics.

The development reduces irrationality to exact non-integrality certificates.
For \(t\in\mathbb{N}\), let
\(H_t=\operatorname{lcm}(1,\ldots,t)\), with \(H_0=1\), and write
\[
\mathrm{Cert}(t)\quad\Longleftrightarrow\quad
\exists L\in\mathbb{N},\ \mathrm{certifiedKill}(H_t,H_t,L).
\]
The \repolink{lean/Erdos249257/TotientTailPeriodKiller.lean\#L70-L74}{certificate predicate}
uses a finite integer truncation of the difference of two totient tails.
It requires the truncation's residue modulo \(2^L\) to avoid both ends of
the residue interval by an explicit tail-error bound. Thus a certificate
can be checked using finite integer arithmetic, although its consequence
concerns an infinite series.

Lean checks the finite theorem
\[
\forall t\le82,\quad \mathrm{Cert}(t).
\]
This is the
\repolink{lean/ErdosProblems/Skip/LadderT67.lean\#L71264-L71271}{finite diagonal-certificate theorem}.
The endpoint needs an unbounded supply:
\[
\forall T,\quad \exists t>T,\quad \mathrm{Cert}(t).
\]
The
\repolink{lean/Erdos249257/LcmConeFlatness.lean\#L426-L435}{diagonal-supply equivalence}
proves that this unbounded statement is equivalent to irrationality.
The source uses \(\forall t_0,\exists t\ge t_0\); over natural numbers this
is equivalent to the strict-cutoff form displayed here.
It does not prove the missing implication from the finite
range to the unbounded statement. No matter how large a fixed checked bound is,
a larger cutoff exists.

The example connects the layers directly. Computation constructs the finite
data that Lean checks through the certificate predicate and finite theorem.
The formal graph records dependencies and the semantic graph describes the
result's role. The claim record fixes the finite range and the open unbounded
requirement, whose quantifier difference the paper explains. Comparator
checks selected exact interfaces; local review selection may rank the family;
the release checker requires the limitation to remain visible.

A historical README edit exposed a gap in the last check. It replaced a
clause saying that the finite cases did \emph{not} supply the open requirement
with one saying that they completed it. Lean was unchanged. The release
checker passed because the prose relationship was unregistered; after it was
added to the claim record, the checker rejected a deliberately false copy.
This demonstrates one failure and repair, without measuring detection rate,
coverage of claim-bearing prose, or mathematical review quality.

Nine of the ten edits were rejected. One escaped because the relationship
had not been registered. The original run logs were not retained. The edits
were authored by the checker's author. Only the escaped edit was reconstructed
after the repair. The other nine edits were not rerun against the extended
checklist. The finite theorem makes no \(t=83\) or cofinal claim.
The \repolink{docs/publication_evidence.json}{historical evidence record}
documents this coverage boundary. The
\repolink{research/experiments/publication_mutations.json}{reconstruction manifest}
and \repolink{scripts/run_publication_mutations.py}{mutation harness}
make the registered edits inspectable but do not replace the missing original
logs. The post-repair witness accepts the current README and rejects a test
copy containing the false clause.

Other lanes illustrate different claim boundaries. Problem 257 has a theorem
for full-support representations; the stronger arbitrary-support statement
remains open. Problem 68 has an exact carry-based equivalence without a
theorem producing the required carries. Problem 251 gives an exact series
reformulation without an irrationality proof, and Problem 243 gives a conditional recovery theorem whose
premises remain in the public claim. The 1041 lane also records a
counterexample and corrected statement. In each case, publication preserves
the quantifiers, premises, and nearest stronger open statement.

\section{Inspection routes}\label{sec:routes}

The following routes lead from common inspection questions to the relevant
public sources.

\begin{table}[H]
\centering\small
\begin{tabular}{@{}>{\raggedright\arraybackslash}p{0.34\linewidth}
                    >{\raggedright\arraybackslash}p{0.54\linewidth}@{}}
\toprule
\papertablehead
Question & Start here \\
\midrule
How does the public repository fit together? &
\repolink{docs/ARCHITECTURE.md}{docs/ARCHITECTURE.md} \\
What may the project say, and what remains open? &
\repolink{docs/claims.json}{docs/claims.json} and
\repolink{docs/methodology.json}{docs/methodology.json} \\
Where is the checked mathematics? &
\repolink{lean/Erdos249257.lean}{lean/Erdos249257.lean} and
\repolink{lean/ErdosProblems.lean}{lean/ErdosProblems.lean} \\
How can I navigate without reading the whole corpus? &
\repolink{docs/ORIENTATION.md}{docs/ORIENTATION.md} and
\repolink{scripts/query_corpus.py}{scripts/query\_corpus.py} \\
What exactly does Comparator check? &
\repolink{docs/EXTERNAL_VERIFICATION.md}{docs/EXTERNAL\_VERIFICATION.md} and
\repolink{verification/comparator.json}{verification/comparator.json} \\
What does the local review selection qualify? &
\repolink{docs/verification/PALOMAR_QUALIFICATION.md}{PALOMAR\_QUALIFICATION.md} and
\repolink{docs/PALOMAR_RESULT_SHOWCASE.json}{docs/PALOMAR\_RESULT\_SHOWCASE.json} \\
Which papers exist and what question does each answer? &
\repolink{docs/papers/README.md}{docs/papers/README.md} \\
Which checks gate a release? &
\repolink{scripts/check_release.py}{scripts/check\_release.py} and
\repolink{.github/workflows/lean.yml}{.github/workflows/lean.yml} \\
How can work return with public credit? &
\repolink{CONTRIBUTING.md}{CONTRIBUTING.md} and
\repolink{docs/research-commons/CONTRIBUTIONS.md}{research-commons/CONTRIBUTIONS.md} \\
\bottomrule
\end{tabular}
\caption{Compact public inspection routes. These are entry points, not a new
authority layer.}
\label{tab:routes}
\end{table}

\section{What can be trusted}\label{sec:trust}
\papersectiontarget{systems-trust}

The architecture supports a chain of narrow conclusions:

\begin{enumerate}
\item a recorded experiment ran on its stated finite inputs;
\item a pinned Lean kernel accepted its stated formal source;
\item a maintainer approved a selected interpretation and public boundary;
\item Comparator matched selected proof-bearing declarations to separately
      stated interfaces and an axiom budget;
\item the local review-selection checks validated the structure of a selected
      review packet;
\item the release program found every configured public relationship intact on
      the named revision.
\end{enumerate}

Together, the receipts trace a result through the system. They do not establish
that its formalisation is uniquely right, its review independent, its
contribution novel or significant, or its registered boundary complete. They
also do not establish a solution to an open Erd\H{o}s problem. A coordinated
mistake in source, record, and prose can satisfy every structural comparison.
Unregistered prose can overclaim, and literature status can change.

Multiple sources, explicit boundaries, and negative tests near high-risk
claims help expose these errors. Mathematical judgement remains necessary,
and the system does not require a second independent mathematician. Assurance
cases have a similar limit: their notation can record a claimed support
relationship without proving that the evidence is true or sufficient
\cite{gsn}.

A capability described in this paper may have any of four kinds of support:

\begin{center}
\small
\begin{tabular}{@{}>{\raggedright\arraybackslash}p{0.20\linewidth}
                    >{\raggedright\arraybackslash}p{0.35\linewidth}
                    >{\raggedright\arraybackslash}p{0.37\linewidth}@{}}
\toprule
Kind of support & Example & What may be said \\
\midrule
Executable check & A session or return validated against required fields and source identity; a Comparator interface match & The program checks the stated conditions on the stated inputs \\
Agent workflow instruction & The discovery and stewardship responsibilities in a public skill & The workflow directs the agent; it does not show that every agent complied \\
Mathematical judgement & Whether a result expresses the intended mathematics or improves on prior work & A named assessment is required; agreement between source, record, and prose is not one \\
External outcome & Another researcher reproduces, adopts, corrects, or builds on the work & A recorded external event is required; the architecture implies none \\
\bottomrule
\end{tabular}
\end{center}

\section{Scaling from one clone to a search network}\label{sec:scaling}
\papersectiontarget{systems-scaling}

\subsection{Clone, attach compute, and mine bounded results}

A fresh public clone fixes a Lean toolchain, exposes supported roots, records
reviewed claims and open boundaries, and supplies query and release programs.
An external researcher can attach an agent runner to these interfaces, assign
disjoint theorem or experiment lanes, and submit candidate changes through the
proof and publication gates. Added compute can increase the number and
diversity of attempts without changing what a passing receipt means.

In the proposed ``result mining'' mode, workers select bounded frontier
objects and return proof attempts, counterexamples, computations, or
literature findings as exact artefacts. Fan-in checks Lean, reconciles
statements, records negative evidence, and ranks the strongest surviving
result families. Each result packet carries provenance, scope, formal status,
mathematical role, and a nearest stronger open statement. Local selection
allocates review attention to these packets, while Comparator protects
selected exact interfaces.

At scale, stewardship would handle each stable delta as it arrives. It
compares the delta with the candidate universe, updates the paper hierarchy
and assurance interfaces when warranted, and returns a ranked frontier to the
workers. Several miners can share one stewardship goal, and one steward can
cover several problems, provided every write and mathematical object retains
one owner. A large run may produce no new mathematical family; a short
negative result may deserve more prominence than many formal lemmas.

The controller admits work according to path conflicts, memory, CPU, disk
pressure, and the expected value of the lane. Many agents can in principle
search disjoint mathematical neighbourhoods, while practical limits differ by
activity. Reading and cheap computation can run widely; changes require path
leases, and heavy Lean validation has narrower capacity.

Public Lean validation uses
\repolink{scripts/validation_singleflight.py}{single-flight request coordination}:
equivalent requests share a detached owner and may reuse its completed receipt.
The request key conservatively includes all visible Lean source content,
build authorities, toolchain and dependency-lock identities, and the normalised
command, which includes the targets. It is independent of the checkout path;
unrelated paper edits do not change the Lean request key. This is content-based
reuse, not a claim that arbitrary equivalent commands or proofs are recognised.

Distinct requests still compete for the same host-wide Mathlib resource.
The worker attempts a nonblocking host lock. If another Lean validation owns
it, the request ends with a typed resource deferral, exit~75, and must be
submitted again after that owner finishes. The transient state named
\texttt{queued} does not promise automatic eventual execution. The ordinary
\repolink{scripts/lean_fast_build.py}{focused-build wrapper} submits and waits
for its result; the explicit submit interface returns a receipt that can be
collected later. This lets a controller advance independent work while an
admitted build runs. A deferral reports resource availability, not a failed
mathematical theorem.

Scaling therefore has four limits: logical search width, safe concurrent
mutation, validation throughput, and expert review. Path leases govern
mutation, admission controls validation, and result ranking directs reviewer
attention. Fan-in integrates the outputs once the required receipts exist.

The return path is already public. A person can fork or clone the repository,
work from a named public commit, and open a pull request or research-progress
issue. A structured return can preserve the contributor, collaborators, tool
operator, disclosed model systems, starting commit, evidence, affected result,
and surviving limitation as separate fields. Only an accepted receipt enters
the generated contribution views. Acceptance, mathematical claim status, and
release inclusion remain separate decisions, so a credited counterexample or
failed route need not be mislabelled as a theorem. Later corrections append a
new history rather than erasing the earlier contributor. Attribution and pull
requests are standard practice; the architectural role here is to carry credit
and provenance through fan-in without allowing either to strengthen the claim.

The named starting commit remains useful even when main advances. Git retains
the common ancestor needed to inspect the contributor's original delta. A
maintainer can replay that state, reconcile the reviewed substance with the
current tree, and rerun current validation and consequence propagation. A
material conflict resolution is a new integration contribution; it does not
rewrite the origin of the earlier work.

The exact contributor-facing sequence is specified in the
\papersectionlink{open-source-mathematics-strategy.pdf}%
{strategy-protocol}{companion contribution protocol}; its credit model is
separated in the
\papersectionlink{open-source-mathematics-strategy.pdf}%
{strategy-credit}{provenance and credit section}. The navigation assumptions
made here are tested by the
\papersectionlink{cold-clone-to-proof-receipt.pdf}%
{cold-clone-problem}{cold-clone case study}.

The public clone supplies the mathematical corpus and its gates. The full
private orchestration layer is not distributed as a turnkey service, and
mass multi-provider mining remains a design target without a reported
benchmark. Public interfaces allow a laboratory to use another scheduler or
producer while retaining the evidence boundary.

No outside contributor had completed this path by 31 August 2026. The
cross-paper links, clone-local skills, and pull-request route are therefore
implemented prototype interfaces whose usability remains an external test.

\subsection{Using the no-go graph}

The proposed graph records endpoint problems, conjectures and reformulations,
mechanisms, computations, formal lemmas, counterexamples, no-go theorems,
reviewed claims, and unresolved obligations. Its edges distinguish implication,
equivalence, dependency, refutation, obstruction of a strategy, weakening,
generalisation, experimental support, formalisation, and publication. A
no-go rules out a class of mechanisms while leaving nearby variants to be
examined separately.

As this negative and positive graph grows, several new uses become testable.
A model may navigate around already closed strategy families, identify a small
cut of missing premises between the current corpus and an endpoint, or search
for analogies between obstruction patterns in different problems. Training
data can be formed from proof/counterexample/no-go triples or from contrastive
pairs consisting of a tempting argument and the exact theorem explaining its
failure. Graph-conditioned models might propose the next useful lemma from a
frontier neighbourhood rather than from the theorem statement alone. These
are hypotheses about future systems, not results established by this paper.

Negative structure also creates risks. An over-broad obstruction edge can hide
a viable variant; repeated failed attempts can teach stylistic avoidance rather
than mathematics; and model-generated semantic edges can disagree with the
formal graph. Training snapshots therefore need immutable generations,
source-level provenance, explicit edge authority, and held-out evaluation.
Lean can verify formal nodes and some edges, but expert judgement remains
necessary for semantic roles and for deciding whether the graph captures the
important mathematical space.

\subsection{An experimental agenda}

Evaluation should compare agents with and without graph context on frontier
selection, repeated dead ends, time to a useful obstruction, premise
discovery, proof completion, and calibration of public claims. It should test
transfer between unrelated mathematical domains as well as neighbouring
Erd\H{o}s problems. Independent teams should replay result packets and assess
the scope of no-go edges. The question is whether the accumulated graph
improves the efficiency and originality of later reasoning, including on
mathematics beyond the corpus already mapped.

\section{Relation to other approaches}\label{sec:related}
\small

\paragraph{Theorem-proving agents.}
LeanDojo and Pantograph provide programmatic Lean environments and proof-state
interaction \cite{leandojo,pantograph}. OpenProver combines a planner,
parallel workers, independent verifiers, compact working memory, a larger
repository, and Lean feedback \cite{openprover}. Agent Hunt studies concurrent
formalisation with locks, bounties, guarded ownership, and collaborative
agents \cite{agenthunt}. DreamProver learns a compact reusable lemma library
through wake--sleep cycles \cite{dreamprover}. These systems provide baselines
for proof interaction, parallel search, and learned proof memory. The present
architecture contributes no new proof-search algorithm; it addresses the
subsequent transition from a found proof to a reviewed public claim.

\paragraph{Graphs and checked exposition.}
Proof blueprints connect informal proof plans to named Lean declarations.
\texttt{leanblueprint} checks that author-supplied declaration names exist,
while LeanArchitect infers formal dependencies and unfinished-proof status and
exports synchronised blueprint material \cite{leanblueprint,leanarchitect}.
The graph layers here have a related navigational role, but the public-claim
record begins after a result has been selected and asks which wording and open
boundary were reviewed.

Semantic-audit systems address a neighbouring problem. Lean Atlas narrows the
declarations a person must inspect for chosen theorem statements, conditional
on the semantic correctness of the returned set and trusted base
\cite{leanatlas}. EconCSLib uses models to translate and compare formal and
informal statements, with saved human judgements for a subset \cite{econcs}.
The present architecture uses models throughout production but assigns none of
them final semantic authority. Their output becomes evidence or a candidate
until a source-specific gate accepts it.

Isabelle/DOF places formal and informal material in a typed checked document
\cite{isadof}; requirements traceability follows commitments across
development and revision \cite{gotel}. This repository leaves prose
unrestricted and records selected public boundaries separately, so
unregistered prose remains outside the checker.

\paragraph{Auditable scientific agents.}
HEP records hypotheses, evidence, belief updates, lineage, and resolution
states in an append-only registry \cite{hep}. Symposium proposes immutable
community publication histories with attributable artefacts, declared
evidence, assumptions, and purpose-sensitive arguments \cite{symposium}.
EurekAgent includes permissions, artefacts, budgets, and human supervision in
its agent environment \cite{eurekagent}. These precedents establish the role
of persistent records and environment design. The comparison here concerns
the separation of experimental evidence, kernel acceptance, reviewed
interpretation, exact-statement assurance, editorial selection, and public
release.

Mutation testing asks whether a test distinguishes a seeded fault from the
original program \cite{demillo,jiaharman}. The deliberately false boundary in
the worked example follows that idea. Because the examples were selected by
the system's author and were not run as a controlled external evaluation, they
show that particular checks can fail; they do not measure overall adequacy.

\paragraph{Contribution and limits.}
The contribution combines existing components into an architecture joining a
private workbench and a self-contained public Lean repository. Formal
dependencies, semantic interpretation, reviewed claims, Comparator
interfaces, review selection, and release relationships have separate
records. Failure modes and formal no-gos pass through the publication process
with positive theorems, while retaining the stronger endpoint that remains
open. Failure logging and hypothesis refutation have precedents; the
comparison here concerns their integration with formal obstruction results,
claim ceilings, and publication assurance. Evidence comes from the design
and a worked reconstruction in one evolving corpus. It provides neither a
controlled performance comparison nor an independent audit or general
reliability estimate.

\paragraph{Scaling beyond this corpus.}
The design is not tied to Erd\H{o}s problems. A new domain needs stable result
identities, source and status records, a formal checker where one exists, and
an explicit public-claim boundary; it need not use the same mathematics or
even Lean. Larger deployments should federate independently owned corpora
rather than build one authority database, separate producer and reviewer
teams, benchmark each boundary independently, and preserve immutable result
generations. Proof-search layers can adopt distributed task markets,
persistent lemma learning, and richer human steering, while publication layers
can add external replay, signed review receipts, and cross-project claim
links. None of these extensions should allow learned process memory or a graph
edge to change mathematical truth by itself.
\normalsize

\section{Conclusion}\label{sec:conclusion}

This architecture organises AI-assisted research around bounded
neighbourhoods in problem-sized mathematical worlds. At each transition from
search to public wording, it records the evidence, the person or program
authorised to assess it, and the stronger statement that remains open.
Reasoning scope, mutation permission, validator capacity, evidence class,
public-claim permission, and reviewer attention have separate gates.

The private workbench supports durable research and concurrent changes.
Type A agents work on claimed paths under checks; Type B agents return
bounded reasoning from web chat. Experiments help choose routes, and Lean
checks formal statements. Graphs support navigation, Comparator checks
selected interfaces, and local selection prepares results for review. The
external Palomar registry checks and filters submissions. Source, claims,
papers, and release receipts are published in a self-contained clone.

The worked example shows how these records preserve a finite theorem's
boundary through publication, including the repair of one escaped overclaim.
It does not establish an automatic mathematician or a verifier for arbitrary
publication claims. Intended meaning still requires mathematical judgement;
independent review, community understanding, and acceptance require evidence
beyond the producing system.

\appendix
\small
\section{Reproducibility}\label{app:repro}

\paragraph{Public first contact.}
A fresh public clone can reproduce the control card and structural checks:
\begin{verbatim}
python3 scripts/proof_cockpit.py --format card
python3 scripts/proof_cockpit.py --check
\end{verbatim}
The first reads committed public metadata; the second checks the claim
registry, cold-clone contract, and orientation freshness.  Neither checks a
proof.  Formal authority begins with the pinned Lean build named by the card.

\paragraph{Source versions and coverage.}
Repository source links in this revision resolve at public commit
\texttt{9b654f4cce44}; they identify inspectable source, not a new validation
run. The \repolink{docs/publication_entry_packet.json}{publication entry packet}
separately identifies the formal-source checkpoint, release identity, and
evidence owners. A reader working from another revision should query those
identities again before comparing its output with this paper.

The \repolink{scripts/query_semantic.py}{semantic query tool}'s
\texttt{coverage} command distinguishes exact proposition evidence, authored
family context, and structural-only linkage. Contextual membership does not
mean that every member received statement-level mathematical review. Coverage
counts belong to the generated receipt at its source digest; they do not
measure review quality, mathematical value, or public-claim completeness.

The \repolink{docs/publication_contract.json}{paper inventory} records source
and PDF cryptographic hashes and validation commands. The
\repolink{docs/publication_evidence.json}{worked-example evidence} and
\repolink{research/experiments/publication_mutations.json}{reconstruction manifest}
specify the deliberately false edits and their limitations. Comparator's
configuration and the local qualification and selection records are linked
in Section~\ref{sec:assurance}.
These artefacts identify what was checked. They do not interpret unrestricted
prose or confer external acceptance.

The private system is described here at the architectural level because it is
production provenance, not a dependency of the public result. Private paths,
operator material, unreleased work, and private ledgers are neither required
nor granted authority by this paper. The public checkout remains the inspection
and replay boundary.

{\tiny
\begin{multicols}{2}
\begin{thebibliography}{20}
\setlength{\itemsep}{0pt}
\setlength{\parsep}{0pt}
\bibitem{lean4} L.~de Moura and S.~Ullrich, \emph{The Lean 4 Theorem Prover
and Programming Language}, in \emph{Automated Deduction---CADE 28}, Lecture
Notes in Computer Science 12699, 2021, pp.~625--635,
\href{https://doi.org/10.1007/978-3-030-79876-5_37}{DOI}.
\bibitem{erdosgraham} P.~Erd\H{o}s and R.~L.~Graham, \emph{Old and New
Problems and Results in Combinatorial Number Theory}, Monographies de
L'Enseignement Math\'ematique 28, 1980, p.~61.
\bibitem{leanblueprint} P.~Massot, \emph{leanblueprint}, plasTeX plugin for
Lean formalisation blueprints, 2020,
\href{https://github.com/PatrickMassot/leanblueprint}{software repository}.
\bibitem{leanarchitect} T.~Zhu, P.~Monticone, S.~Welleck, and J.~Avigad,
\emph{LeanArchitect: Automating Blueprint Generation for Humans and AI},
in \emph{17th International Conference on Interactive Theorem Proving},
LIPIcs 382, 2026, pp.~25:1--25:16.
\bibitem{leanatlas} B.~Yanahama and A.~Sannai, \emph{Lean Atlas: An
Integrated Proof Environment for Scalable Human--AI Collaborative
Formalization}, 2026, \href{https://doi.org/10.48550/arXiv.2604.16347}{arXiv}.
\bibitem{econcs} N.~Garg, \emph{EconCSLib: AI-Assisted Lean Formalization for
Economics \& Computation Research}, 2026,
\href{https://doi.org/10.48550/arXiv.2606.13306}{arXiv}.
\bibitem{isadof} A.~D.~Brucker and B.~Wolff, \emph{Isabelle/DOF: Design and
Implementation}, in \emph{Software Engineering and Formal Methods}, 2019,
pp.~275--293.
\bibitem{gotel} O.~C.~Z.~Gotel and A.~C.~W.~Finkelstein, \emph{An analysis
of the requirements traceability problem}, Proc. First IEEE International
Conference on Requirements Engineering, 1994, pp.~94--101.
\bibitem{demillo} R.~A.~DeMillo, R.~J.~Lipton, and F.~G.~Sayward, \emph{Hints
on test data selection}, IEEE Computer 11(4), 1978, pp.~34--41.
\bibitem{jiaharman} Y.~Jia and M.~Harman, \emph{An analysis and survey of the
development of mutation testing}, IEEE Transactions on Software Engineering
37(5), 2011, pp.~649--678.
\bibitem{gsn} SCSC Assurance Case Working Group, \emph{Goal Structuring
Notation Community Standard, Version 3}, SCSC-141C, May 2021.
\bibitem{taoai} T.~Tao, \emph{Mathematics in the age of AI}, preprint, 2026,
\href{https://doi.org/10.48550/arXiv.2608.16753}{arXiv}.
\bibitem{mathlib} Lean community, \emph{Contributing to mathlib},
\href{https://leanprover-community.github.io/contribute/index.html}%
{contributor guide}, accessed August 2026.
\bibitem{leandojo} K.~Yang et al., \emph{LeanDojo: Theorem Proving with
Retrieval-Augmented Language Models}, NeurIPS 2023.
\bibitem{pantograph} J.~Storrs et al., \emph{Pantograph: A Machine-to-Machine
Interaction Interface for Advanced Theorem Proving, High Level Reasoning, and
Data Extraction in Lean 4}, 2024,
\href{https://doi.org/10.48550/arXiv.2410.16429}{arXiv}.
\bibitem{openprover} M.~Kripner and M.~Straka, \emph{OpenProver: Agentic and
Interactive Theorem Proving with Lean 4}, 2026,
\href{https://doi.org/10.48550/arXiv.2607.09217}{arXiv}.
\bibitem{agenthunt} C.~E.~Brown, C.~Kaliszyk, and J.~Urban, \emph{Agent Hunt:
Bounty Based Collaborative Autoformalization With LLM Agents}, 2026,
\href{https://doi.org/10.48550/arXiv.2603.06737}{arXiv}.
\bibitem{dreamprover} Y.~Zhang et al., \emph{DreamProver: Evolving
Transferable Lemma Libraries via a Wake--Sleep Theorem-Proving Agent}, 2026,
\href{https://doi.org/10.48550/arXiv.2604.26311}{arXiv}.
\bibitem{hep} I.~Takahara and T.~Mizoguchi, \emph{Toward Auditable AI
Scientists: A Hypothesis Evolution Protocol for LLM Agents}, 2026,
\href{https://doi.org/10.48550/arXiv.2607.09195}{arXiv}.
\bibitem{symposium} D.~Pratt, \emph{Symposium: Trust via Auditable Records for
Communities of AI Scientist Agents}, 2026,
\href{https://doi.org/10.48550/arXiv.2608.19511}{arXiv}.
\bibitem{eurekagent} A.~Xin et al., \emph{EurekAgent: Agent Environment
Engineering is All You Need for Autonomous Scientific Discovery}, 2026,
\href{https://doi.org/10.48550/arXiv.2606.13662}{arXiv}.
\end{thebibliography}
\end{multicols}
}
\end{document}
